\pdfoutput=1
\documentclass[indonesian]{shinybook}
\usepackage{tikz}
\usepackage{pgfplots}
\usepgfplotslibrary{groupplots}
\pgfplotsset{compat=1.18}
\usetikzlibrary{intersections,calc}
\usetikzlibrary{arrows.meta, decorations.markings}
\usepackage{glossaries-extra}
\usepackage{needspace}

\newcommand{\half}{\frac{1}{2}}
\newcommand{\N}{\mathbb{N}}
\newcommand{\Z}{\mathbb{Z}}
\newcommand{\Q}{\mathbb{Q}}
\newcommand{\R}{\mathbb{R}}
\newcommand{\Rb}{\overline{\mathbb{R}}}
\newcommand{\bR}{\overline{\mathbb{R}}}
\newcommand{\F}{\mathbb{F}}
\newcommand{\1}{\mathbb{1}}
\newcommand{\linop}{\mathcal{L}}
\newcommand{\Rnn}{\mathbb{R}^{n\times n}}
\newcommand{\Cnn}{\mathbb{C}^{n\times n}}
\newcommand{\E}{\mathbb{E}}
\usepackage{pgffor}
\foreach \x in {A,...,Z}{%
    \expandafter\xdef\csname cal\x\endcsname{\noexpand\ensuremath{\noexpand\mathcal{\x}}}
}

\renewcommand{\implies}{\Rightarrow}
\newcommand{\setto}{\rightrightarrows}
\newcommand{\equivalent}{\Leftrightarrow}
\newcommand{\eps}{\varepsilon}
\renewcommand{\phi}{\varphi}
\newcommand{\supp}{\operatorname{\mathrm{supp}}}
\renewcommand{\div}{\operatorname{\mathrm{div}}}
\newcommand{\dom}{\operatorname{\mathrm{dom}}}
\renewcommand{\ker}{\operatorname{\mathrm{ker}}}
\newcommand{\conv}{\operatorname{\mathrm{conv}}}
\newcommand{\rg}{\operatorname{\mathrm{rg}}}
\newcommand{\graph}{\operatorname{\mathrm{graph}}}
\newcommand{\tr}{\operatorname{\mathrm{tr}}}
\newcommand{\epi}{\operatorname{\mathrm{epi}}}
% Never redefine TeX's primitive \span: amsmath alignment environments use it.
% Use \spann for the linear-span operator in reader prose and formulae.
\newcommand{\sign}{\operatorname{\mathrm{sign}}}
\DeclareMathOperator{\Id}{\mathrm{Id}}
\newcommand{\prox}{\mathrm{prox}}
\newcommand{\proj}{\mathrm{proj}}
\newcommand{\spann}{\operatorname{\mathrm{span}}}
\newcommand{\Lloc}{L^1_{\mathrm{loc}}}
\newcommand{\Cinf}{C^\infty_0}

% Nama lingkungan dilokalkan; penomoran dan topologi sumber dipertahankan.
\newtheorem{theorem}{Teorema}[chapter]
\newtheorem{lemma}[theorem]{Lema}
\newtheorem{cor}[theorem]{Korolari}
\theoremstyle{definition}
\newtheorem{defn}[theorem]{Definisi}
\newtheorem{prop}[theorem]{Proposisi}
\newtheorem{example}[theorem]{Contoh}
\newtheorem{assumption}[theorem]{Asumsi}
\newtheorem{exercise}[theorem]{Latihan}
\newtheorem*{rem}{Catatan}

\crefname{theorem}{teorema}{teorema}
\Crefname{theorem}{Teorema}{Teorema}
\crefname{lemma}{lema}{lema}
\Crefname{lemma}{Lema}{Lema}
\crefname{cor}{korolari}{korolari}
\Crefname{cor}{Korolari}{Korolari}
\crefname{prop}{proposisi}{proposisi}
\Crefname{prop}{Proposisi}{Proposisi}
\crefname{defn}{definisi}{definisi}
\Crefname{defn}{Definisi}{Definisi}
\crefname{example}{contoh}{contoh}
\Crefname{example}{Contoh}{Contoh}
\crefname{assumption}{asumsi}{asumsi}
\Crefname{assumption}{Asumsi}{Asumsi}
\crefname{rem}{catatan}{catatan}
\Crefname{rem}{Catatan}{Catatan}
\crefname{exercise}{latihan}{latihan}
\Crefname{exercise}{Latihan}{Latihan}

\renewcommand{\proofname}{Bukti}
\renewcommand{\contentsname}{Daftar Isi}
\renewcommand{\figurename}{Gambar}

\newcommand{\norm}[1]{\|#1\|}
\newcommand{\abs}[1]{\left|#1\right|}
\newcommand{\inner}[2]{\left\langle#1 , #2\right\rangle}
\newcommand{\dual}[1]{\langle #1 \rangle}
\newcommand{\setof}[2]{\left\{#1 \;\middle|\; #2\right\}}
\newcommand{\wkto}{\rightharpoonup}
\DeclareMathOperator{\Exp}{\mathbb{E}}
\DeclareMathOperator{\Var}{\mathbb{V}}
\DeclareMathOperator{\Cov}{\mathrm{Cov}}
\newcommand{\Prob}{\mathbb{P}}

\usepackage{enumerate}
\newcommand{\weaki}{\protect{\u{\i}}}
\AtBeginEnvironment{remark}{\small}

\ifdefined\OIntegratedReader
  % Keep examples as ordinary theorem blocks in the tagged master. The legacy
  % mdframed/tablefootnote patches interfere with the synchronized latex-lab
  % hooks, while no canonical body calls \tablefootnote.
\else
  \usepackage{mdframed}
  \mdfdefinestyle{examplebg}{
      backgroundcolor=structure!7,%
      hidealllines=true,%
      splittopskip=\topskip,
      frametitleaboveskip=0,
      frametitlebelowskip=0,
      innertopmargin=-.3em,%
      innerbottommargin=.7em,%
      innerleftmargin=.7em,%
      innerrightmargin=.7em,%
  }
  \surroundwithmdframed[style=examplebg]{example}
  \usepackage{tablefootnote}
  \makeatletter
  \AfterEndEnvironment{mdframed}{%
   \tfn@tablefootnoteprintout%
   \gdef\tfn@fnt{0}%
  }
  \makeatother
\fi

\usepackage{enumitem}
\setlist[enumerate]{label=(\roman*)}
\usepackage{scrhack}

\newcommand{\ie}{\textit{yaitu}}
\newcommand{\eg}{\textit{misalnya}}
\newcommand{\cf}{\textit{bandingkan}}
\newcommand{\emptyarg}{\,\cdot\,}
\newcommand{\dd}{\mathrm{d}}
\newcommand{\todo}[1]{{\color{red}Perlu dikerjakan: #1}}
\newcommand{\Oc}{\mathcal{O}}
\newcommand{\Lc}{\mathcal{L}}
\newcommand{\Gc}{\mathcal{G}}
\newcommand{\Ic}{\mathcal{I}}
\newcommand{\Pc}{\mathcal{P}}
\newcommand{\Mc}{\mathcal{M}}
\renewcommand{\epsilon}{\varepsilon}
\newcommand{\diag}{\mathrm{diag}}

\microtypesetup{expansion=false}
\pdfmapfile{+libertine.map}
\pdfmapfile{+newtx.map}
\crefname{equation}{persamaan}{persamaan}
\Crefname{equation}{Persamaan}{Persamaan}

\title{Optimisasi Konveks}
\subtitle{Unit 6: Akselerasi\\Edisi Bahasa Indonesia}
\author{Andreas Habring}
\date{Sumber edisi v1: 13 Juli 2026}
\publishers{\small
Terjemahan kerja bahasa Indonesia dari \emph{Lecture Notes: Convex Optimization}, arXiv:2607.11664v1.\\
Sumber asli dan terjemahan ini dilisensikan CC BY 4.0. Perubahan matematis dijelaskan dalam catatan koreksi.\\
Terjemahan mandiri ini bukan karya resmi atau dukungan Andreas Habring maupun TU Graz.}
\hypersetup{
    pdftitle={Optimisasi Konveks - Unit 6: Akselerasi - Edisi Bahasa Indonesia},
    pdfauthor={Andreas Habring},
    pdflang={id-ID},
}

\begin{document}
\maketitle

\frontmatter
\chapter*{Tentang unit ini}
Unit ini menerjemahkan secara utuh Bab 6, \emph{Acceleration}, dari versi sumber yang ditentukan di atas. Unit ini melanjutkan Unit 5. Pembaca diharapkan telah mengenal fungsi konveks dan konveks kuat, gradien Lipschitz, Hessian, radius spektral, operator proksimal, serta metode gradien proksimal.

Urutan hasil, struktur bukti, seluruh rumus, tujuh pengenal rujukan silang, dan ajakan verifikasi sumber dipertahankan. Sejumlah kekeliruan matematis dalam sumber diperbaiki secara terbuka; ringkasannya disertakan pada akhir unit dan rekaman terstruktur lengkap menyertai sumber edisi ini. Tidak ada sitasi, gambar, atau catatan kaki dalam bab sumber. Tidak ada layanan daring atau perangkat lunak tertutup yang diperlukan untuk membaca unit ini.

\tableofcontents
\mainmatter
\setcounter{chapter}{5}
\chapter{Akselerasi}

% segment-id: d90.hab.v1.ch06.seg0001
Untuk setiap \(L>0\), titik awal \(x_0\in\R^d\), bilangan bulat \(k\geq1\), dan \(d\geq2k+1\), serta untuk setiap metode orde pertama yang iterasinya memenuhi
\begin{equation}
    x_k\in x_0+\mathrm{span}\{\nabla f(x_0),\dots,\nabla f(x_{k-1})\}
\end{equation}
terdapat fungsi konveks \(f:\R^d\rightarrow\R\) yang bergradien \(L\)-Lipschitz dan mempunyai peminim \(x^*\), sedemikian sehingga
\begin{equation}
    f(x_k)-f(x^*)\geq\frac{3L\|x_0-x^*\|^2}{32(k+1)^2}.
\end{equation}
Perhatikan bahwa sejauh ini kita belum memperoleh laju konvergensi yang lebih baik daripada \(\Oc(1/k)\). Menutup kesenjangan menuju \(\Oc(1/k^2)\) merupakan tujuan bab ini.

% segment-id: d90.hab.v1.ch06.seg0002
\section{Metode bola berat Polyak}

Penurunan gradien untuk meminimumkan masalah tak berkendala
\begin{equation}
    \min_x f(x)
\end{equation}
dapat ditafsirkan sebagai diskretisasi dari \emph{aliran gradien}
\begin{equation}
    \dot x=-\nabla f(x).
\end{equation}
Jika fungsi \(f\) ditafsirkan sebagai potensial, maka \(-\nabla f\) merupakan gaya yang ditimbulkan potensial tersebut pada sebuah partikel bermassa satu. Secara alami kita dapat menambahkan gesekan ke dalam model ini. Gaya gesek biasanya dimodelkan sebanding dengan kecepatan, sehingga diperoleh
\begin{equation}\label{acceleration:eq:friction_ode}
    \ddot x=-\gamma\dot x-\nabla f(x).
\end{equation}
Kita mendiskretkan persamaan diferensial biasa ini melalui \(\dot x(t)\approx\frac{x(t)-x(t-h)}{h}\) dan
\begin{equation}
    \ddot x(t)
    \approx\frac{\dot x(t+h)-\dot x(t)}{h}
    \approx\frac{x(t+h)-2x(t)+x(t-h)}{h^2}.
\end{equation}
Dengan menamai \(x(t-h)=x_{k-1}\), \(x(t)=x_k\), dan \(x(t+h)=x_{k+1}\), substitusi ke dalam~\cref{acceleration:eq:friction_ode} memberikan
\begin{equation}
    x_{k+1}-2x_k+x_{k-1}
    =-h\gamma(x_k-x_{k-1})-h^2\nabla f(x_k).
\end{equation}
Setelah ditata ulang, pembaruannya menjadi
\begin{equation}
    x_{k+1}=x_k+(1-\gamma h)(x_k-x_{k-1})-h^2\nabla f(x_k).
\end{equation}
Dengan menamai ulang konstanta dan mengizinkannya bergantung pada iterasi, kita sampai pada \emph{metode bola berat Polyak}
\begin{equation}
    x_{k+1}=x_k+\beta_k(x_k-x_{k-1})-\tau_k\nabla f(x_k).
\end{equation}
Bahkan tanpa penurunan dari sudut pandang persamaan diferensial, aturan pembaruan ini cukup intuitif: selain arah penurunan bagi \(f\), pembaruan tersebut memuat suku \emph{inersia} yang menambahkan komponen searah dengan langkah sebelumnya.

% segment-id: d90.hab.v1.ch06.seg0003
Bukti konvergensi metode bola berat akan memakai hasil dasar aljabar linear berikut, yang disertakan demi kelengkapan.

\begin{lemma}\label{acceleration:lemma:spectral_radius}
    Untuk setiap \(A\in\R^{d\times d}\), dengan \(\|\cdot\|\) norma operator yang diinduksi oleh norma Euclidean, berlaku \(\lim_{n\rightarrow\infty}\|A^n\|^{1/n}=\rho(A)\), dengan
    \begin{equation}
        \rho(A)=\max_i|\lambda_i(A)|
    \end{equation}
    dan \(\lambda_i(A)\) nilai eigen \(A\). Kesimpulan yang sama berlaku bagi setiap norma matriks yang kompatibel, karena semua norma pada ruang matriks berdimensi hingga ekuivalen.
\end{lemma}
\begin{proof}
    Kompleksifikasikan \(A\) ke operator pada \(\mathbb{C}^d\). Norma operator Euclidean untuk matriks real sama pada \(\R^d\) dan \(\mathbb{C}^d\). Ambil nilai eigen kompleks \(\lambda\) dengan \(|\lambda|=\rho(A)\) serta vektor eigennya \(v\in\mathbb{C}^d\) yang dinormalisasi. Maka
    \begin{equation}
        \|A^n\|^{1/n}\geq\|A^nv\|^{1/n}=|\lambda|=\rho(A),
    \end{equation}
    sehingga \(\liminf_{n\rightarrow\infty}\|A^n\|^{1/n}\geq\rho(A)\).

    Untuk ketaksamaan sebaliknya, gunakan bentuk normal Jordan kompleks
    \begin{equation}
        A=BDB^{-1}.
    \end{equation}
    Setelah blok-blok Jordan diurutkan, matriks \(D\) dapat ditulis sebagai
    \begin{equation}
        D=\begin{pmatrix}
            \lambda_1&\varepsilon_1&\dots&\dots&\\
            0&\lambda_2&\varepsilon_2&\dots&\\
            &&\ddots&\ddots&\\
            &&\dots&\lambda_{d-1}&\varepsilon_{d-1}\\
            &&&0&\lambda_d
        \end{pmatrix},
    \end{equation}
    dengan \(\varepsilon_i\in\{0,1\}\), dan \(\varepsilon_i=1\) hanya jika posisi \(i\) dan \(i+1\) berada dalam blok Jordan yang sama sehingga \(\lambda_i=\lambda_{i+1}\). Tuliskan \(D=\Lambda+N\), dengan \(\Lambda\) diagonal dan \(N\) bagian nilpoten. Kedua matriks ini komutatif. Jika ukuran blok Jordan terbesar adalah \(m\), maka \(N^m=0\). Untuk \(\rho(A)>0\) dan \(n\geq m-1\), ekspansi binomial memberi
    \begin{equation}
        \begin{aligned}
            \|D^n\|
            &=\left\|\sum_{j=0}^{m-1}
            \begin{pmatrix}
                n\\
                j
            \end{pmatrix}\Lambda^{n-j}N^j\right\|\\
            &\leq\sum_{j=0}^{m-1}
            \begin{pmatrix}
                n\\
                j
            \end{pmatrix}\|\Lambda^{n-j}\|\,\|N^j\|\\
            &\leq C_{\rho}\,n^{m-1}\rho(A)^{\,n-m+1},
        \end{aligned}
    \end{equation}
    untuk suatu konstanta \(C_{\rho}>0\) yang tidak bergantung pada \(n\). Jika \(\rho(A)=0\), semua elemen diagonal nol, sehingga \(D=N\) dan \(D^n=0\) bagi \(n\geq m\); jadi kesimpulan langsung berlaku. Untuk \(\rho(A)>0\), batas di atas menghasilkan
    \begin{equation}
        \begin{aligned}
            \|D^n\|^{1/n}
            \leq C_{\rho}^{1/n}n^{(m-1)/n}
            \rho(A)^{\,1-(m-1)/n}.
        \end{aligned}
    \end{equation}
    Karena \(C_{\rho}^{1/n}\rightarrow1\), \(n^{1/n}\rightarrow1\), dan \(\rho(A)^{-(m-1)/n}\rightarrow1\), kita memperoleh
    \begin{equation}
        \begin{aligned}
            \limsup_{n\rightarrow\infty}\|D^n\|^{1/n}
            \leq\rho(A).
        \end{aligned}
    \end{equation}
    Akhirnya,
    \begin{equation}
        \limsup_{n\rightarrow\infty}\|A^n\|^{1/n}
        =\limsup_{n\rightarrow\infty}\|BD^nB^{-1}\|^{1/n}
        \leq\limsup_{n\rightarrow\infty}
        \|B\|^{1/n}\|D^n\|^{1/n}\|B^{-1}\|^{1/n}
        \leq\rho(A),
    \end{equation}
    yang, bersama batas bawah, membuktikan hasilnya.
\end{proof}

% segment-id: d90.hab.v1.ch06.seg0004
\begin{cor}\label{acceleration:cor:spectral_radius}
    Berlaku \(\lim_{k\rightarrow\infty}A^k=0\) jika dan hanya jika \(\rho(A)<1\). Selain itu, dalam hal ini, untuk setiap \(\epsilon>0\) terdapat \(c(\epsilon)>0\) sedemikian sehingga \(\|A^k\|\leq c(\epsilon)(\rho(A)+\epsilon)^k\).
\end{cor}
\begin{proof}
    Jika \(A^k\rightarrow0\), maka untuk setiap pasangan eigen kompleks \((\lambda,v)\) berlaku \(A^kv=\lambda^kv\rightarrow0\); akibatnya \(|\lambda|<1\) dan \(\rho(A)<1\). Sebaliknya, jika \(\rho(A)<1\), pilih \(0<\epsilon_0<1-\rho(A)\). Berdasarkan~\cref{acceleration:lemma:spectral_radius}, terdapat \(n\) sedemikian sehingga untuk \(k\geq n\),
    \begin{equation}
        \|A^k\|\leq(\rho(A)+\epsilon_0)^k,
    \end{equation}
    yang menuju nol. Untuk pernyataan kuantitatif dan sembarang \(\epsilon>0\), \cref{acceleration:lemma:spectral_radius} memberikan batas tersebut untuk semua \(k\) yang cukup besar. Dengan memperbesar konstanta untuk indeks awal yang berhingga banyaknya, kita dapat memilih
    \begin{equation}
        c(\epsilon)\geq
        \max_{0\leq k<n}
        \frac{\|A^k\|}{(\rho(A)+\epsilon)^k},
    \end{equation}
    sehingga batas berlaku bagi setiap \(k\geq0\).
\end{proof}

% segment-id: d90.hab.v1.ch06.seg0005
Sekarang kita dapat membuktikan konvergensi metode bola berat. Tekniknya baku: kita menganalisis nilai eigen linearisasi aturan pembaruan.

\begin{theorem}[Konvergensi bola berat]
    Misalkan \(f\in C^2(U)\) pada suatu lingkungan \(U\subset\R^d\) dari peminim stasioner lokal \(x^*\), dengan
    \begin{equation}
        \nabla f(x^*)=0,\qquad
        0<\mu\leq L,\qquad
        \mu\Id\preceq\nabla^2f(x^*)\preceq L\Id,
        \qquad 0\leq\beta<1,\quad
        0<\tau<\frac{2(1+\beta)}{L}.
    \end{equation}
    Gunakan parameter konstan \(\beta_k=\beta\) dan \(\tau_k=\tau\). Terdapat \(\varepsilon>0\) sedemikian sehingga, jika \(x_0,x_1\in B_\varepsilon(x^*)\), maka \(x_k\rightarrow x^*\). Lebih khusus, dengan
    \begin{equation}
        q=\rho(\mathcal{H})<1,\qquad
        \|x_k-x^*\|\leq
        c(\delta)(q+\delta)^{k-1}
        \|(x_1-x^*,x_0-x^*)\|
    \end{equation}
    untuk setiap \(0<\delta<1-q\). Jika yang diketahui tentang kurvatur hanya selang \([\mu,L]\), parameter yang meminimumkan radius spektral kasus terburuk dari model kuadratik atau linearisasi lokal adalah
    \begin{equation}
        q_{\mathrm{wc}}=
        \frac{\sqrt{L}-\sqrt{\mu}}{\sqrt{L}+\sqrt{\mu}},
        \qquad
        \tau=\frac{4}{(\sqrt{L}+\sqrt{\mu})^2},
        \qquad
        \beta=
        \left(\frac{\sqrt{L}-\sqrt{\mu}}
        {\sqrt{L}+\sqrt{\mu}}\right)^2.
    \end{equation}
\end{theorem}

% segment-id: d90.hab.v1.ch06.seg0006
\begin{proof}
    Tetapkan \(u_k=x_k-x^*\). Karena \(\nabla f(x^*)=0\), rekursi kedalaman dua dapat ditulis tepat sebagai
    \begin{equation}
        \begin{aligned}
            \begin{bmatrix}
                u_{k+1}\\
                u_k
            \end{bmatrix}
            &=
            \begin{bmatrix}
                (1+\beta)u_k-\beta u_{k-1}
                -\tau\bigl(\nabla f(x^*+u_k)-\nabla f(x^*)\bigr)\\
                u_k
            \end{bmatrix}\\
            &=
            \begin{pmatrix}
                (1+\beta)\Id&-\beta\Id\\
                \Id&0
            \end{pmatrix}
            \begin{bmatrix}
                u_k\\
                u_{k-1}
            \end{bmatrix}
            -\tau\bigl(\nabla f(x^*+u_k)-\nabla f(x^*),0\bigr).
        \end{aligned}
    \end{equation}
    Untuk menyiapkan linearisasi, definisikan vektor keadaan dan matriks linear
    \begin{equation}
        \begin{aligned}
            s_k&\coloneqq
            \begin{bmatrix}
                u_k\\
                u_{k-1}
            \end{bmatrix},
            \qquad
            \mathcal{H}&\coloneqq
            \begin{pmatrix}
                (1+\beta)\Id-\tau\nabla^2f(x^*)&-\beta\Id\\
                \Id&0
            \end{pmatrix},
            \qquad
            s_{k+1}\coloneqq
            \begin{bmatrix}
                u_{k+1}\\
                u_k
            \end{bmatrix}.
        \end{aligned}
    \end{equation}
    Teorema dasar kalkulus memberikan
    \begin{equation}
        \begin{aligned}
            \nabla f(x^*+u_k)-\nabla f(x^*)
            &=\int_0^1\nabla^2f(x^*+tu_k)\,u_k\,\dd t\\
            &=\nabla^2f(x^*)u_k+r_k,\qquad
            r_k\coloneqq
            \left[\int_0^1
            \bigl(\nabla^2f(x^*+tu_k)-\nabla^2f(x^*)\bigr)\dd t\right]u_k.
        \end{aligned}
    \end{equation}
    Dengan demikian, rekursi keadaan yang tepat ialah
    \begin{equation}
        \begin{aligned}
            \begin{bmatrix}
                u_{k+1}\\
                u_k
            \end{bmatrix}
            &=
            \underbrace{\begin{pmatrix}
                (1+\beta)\Id-\tau\nabla^2f(x^*)&-\beta\Id\\
                \Id&0
            \end{pmatrix}}_{\mathcal{H}}
            \begin{bmatrix}
                u_k\\
                u_{k-1}
            \end{bmatrix}
            +
            \begin{bmatrix}
                -\tau r_k\\
                0
            \end{bmatrix}.
        \end{aligned}
    \end{equation}
    Kontinuitas Hessian menunjukkan \(r_k=o(\|u_k\|)\) ketika \(u_k\rightarrow0\).

% segment-id: d90.hab.v1.ch06.seg0007
    Berdasarkan~\cref{acceleration:cor:spectral_radius}, kita menganalisis nilai eigen \(\mathcal{H}\). Jika \((v,w)\) merupakan vektor eigen dengan nilai eigen \(\lambda\), maka
    \begin{equation}
        \begin{cases}
            \bigl((1+\beta)\Id-\tau\nabla^2f(x^*)\bigr)v-\beta w
            &=\lambda v,\\
            v&=\lambda w.
        \end{cases}
    \end{equation}
    Dengan menghitung determinan blok, atau mengeliminasi relasi di atas tanpa pernah membagi dengan \(\lambda\), setiap nilai eigen berkaitan dengan suatu \(\eta\in[\mu,L]\) dan memenuhi
    \begin{equation}
        p_\eta(\lambda)
        \coloneqq\lambda^2-(1+\beta-\tau\eta)\lambda+\beta=0.
    \end{equation}
    Jadi, dengan \(a_\eta=1+\beta-\tau\eta\), penyempurnaan kuadrat memberikan
    \begin{equation}\label{acceleration:eq:heavy_ball1}
        \left(\lambda-\frac{a_\eta}{2}\right)^2
        =\left(\frac{a_\eta}{2}\right)^2-\beta.
    \end{equation}
    Akar dalam~\eqref{acceleration:eq:heavy_ball1} berada ketat di dalam lingkaran satuan jika dan hanya jika kriteria Schur--Jury untuk polinom monik kuadrat real berikut terpenuhi:
    \begin{equation}
        |\beta|<1,\qquad
        1-a_\eta+\beta>0,\qquad
        1+a_\eta+\beta>0,
        \quad\Longleftrightarrow\quad
        \beta<1,\quad\tau\eta>0,\quad
        \tau\eta<2(1+\beta).
    \end{equation}
    Hipotesis teorema memenuhi syarat ini bagi semua \(\eta\in[\mu,L]\), termasuk kasus \(\beta=0\); oleh sebab itu \(q=\rho(\mathcal{H})<1\).

    Ambil \(0<\delta<1-q\). Terdapat norma ekuivalen \(\|\cdot\|_\delta\) pada \(\R^{2d}\) yang norma operator terinduksinya memenuhi \(\|\mathcal{H}\|_\delta\leq q+\delta/2\). Karena \(r(u)=o(\|u\|)\) dan semua norma berdimensi hingga ekuivalen, lingkungan \(x^*\) dapat diperkecil sehingga suku gangguan memenuhi batas yang diperlukan. Maka rekursi keadaan menghasilkan
    \begin{equation}
        \|s_{k+1}\|_\delta
        \leq\|\mathcal{H}\|_\delta\|s_k\|_\delta
        +\|(-\tau r_k,0)\|_\delta
        \leq(q+\delta)\|s_k\|_\delta,
        \qquad
        \|s_k\|_\delta
        \leq(q+\delta)^{k-1}\|s_1\|_\delta.
    \end{equation}
    Dengan memilih \(x_0,x_1\) cukup dekat ke \(x^*\), bola kecil yang digunakan di atas tetap invarian. Ekuivalensi norma lalu memberikan estimasi dalam pernyataan teorema dan konvergensi lokal.

    Tinggal membuktikan optimalitas kasus terburuk yang dinyatakan. Kasus \(L=\mu\) langsung: pilihan \(\beta=0\) dan \(\tau=1/L\) membuat \(a_\eta=0\) dan kedua akar nol. Sekarang andaikan \(L>\mu\). Untuk setiap radius kasus terburuk yang diklaim \(r>0\), kasus \(r\geq1\) tidak dapat memperbaiki laju yang akan diperoleh, jadi cukup tinjau \(0<r<1\). Terapkan kriteria Schur--Jury setelah penskalaan \(\lambda=rz\). Agar kedua akar setiap polinom \(p_\eta\) bermodulus paling besar \(r\), perlu
    \(0\leq\beta\leq r^2\) dan
    \(\lvert a_\mu\rvert,\lvert a_L\rvert\leq R\), dengan
    \(R\coloneqq r+\beta/r\).

    Tetapkan \(h=(L-\mu)/(L+\mu)\). Karena
    \(L a_\mu-\mu a_L=(L-\mu)(1+\beta)\), kelayakan kedua titik ujung mengharuskan
    \(R\geq h(1+\beta)\). Selisih
    \(R-h(1+\beta)=r-h+\beta(1/r-h)\) meningkat terhadap \(\beta\) pada \(0\leq\beta\leq r^2\), sebab \(0<r<1\) dan \(0<h<1\). Oleh karena itu, bahkan nilai maksimumnya harus tak negatif:
    \(2r-h(1+r^2)\geq0\). Jadi \(2r\geq h(1+r^2)\), yang pada \(0<r<1\) ekuivalen dengan
    \(r\geq(\sqrt L-\sqrt\mu)/(\sqrt L+\sqrt\mu)\).

    Kesamaan memaksa \(\beta=r^2\) serta
    \(a_\mu=2r\) dan \(a_L=-2r\). Menyelesaikan kedua hubungan ujung tersebut menghasilkan
    \begin{equation}
        \begin{aligned}
            1+\beta-\tau\mu&=2\sqrt{\beta},\qquad
            1+\beta-\tau L=-2\sqrt{\beta}\\
            &\Longrightarrow\quad
            \sqrt{\beta}=q_{\mathrm{wc}}
            =\frac{\sqrt L-\sqrt\mu}{\sqrt L+\sqrt\mu},
            \qquad
            \tau=\frac{4}{(\sqrt L+\sqrt\mu)^2}.
        \end{aligned}
    \end{equation}
    Dengan pilihan ini, \(a_\eta\) bergerak linear dari \(2q_{\mathrm{wc}}\) ke \(-2q_{\mathrm{wc}}\), sehingga semua akar untuk \(\eta\in[\mu,L]\) bermodulus paling besar \(q_{\mathrm{wc}}\). Batas bawah di atas membuktikan bahwa nilai maksimum tersebut tidak dapat diperkecil. Optimalitas ini merupakan optimalitas radius spektral kasus terburuk bagi fungsi kuadratik dan bagi model asimtotik lokal \(C^2\), bukan jaminan konvergensi global bagi semua fungsi konveks mulus.
\end{proof}

\paragraph{Latihan verifikasi.}
Verifikasikan kriteria Schur--Jury di atas langsung dari rumus akar, termasuk kasus \(\beta=0\), dan turunkan kembali parameter kasus terburuk dengan menyamakan perilaku pada \(\eta=\mu\) dan \(\eta=L\).

% segment-id: d90.hab.v1.ch06.seg0008
\section{Akselerasi Nesterov}
Metode Polyak optimal dalam arti radius spektral kasus terburuk bagi fungsi kuadratik konveks kuat, dan secara lokal bagi model \(C^2\) di sekitar peminim. Akselerasi berikut mencapai laju orde optimal di bawah asumsi global yang lebih lemah. Kita kembali meninjau masalah
\begin{equation}
    \min_{x\in\R^d}F(x)\coloneqq f(x)+g(x),
\end{equation}
dengan \(f:\R^d\rightarrow\R\) konveks, terdiferensialkan, dan bergradien \(L\)-Lipschitz untuk \(L>0\), sedangkan \(g:\R^d\rightarrow(-\infty,\infty]\) proper, semikontinu bawah, dan konveks. Andaikan \(\arg\min F\neq\emptyset\). Dengan \(x_0=x_1\), \(t_1=1\), ukuran langkah tetap \(0<\tau\leq1/L\), dan \(k=1,2,\dots\), pembaruan \emph{fast iterative shrinkage-thresholding algorithm} (FISTA, atau metode gradien proksimal cepat) ialah
\begin{equation}
    \begin{cases}
        t_{k+1}&=\dfrac{1+\sqrt{1+4t_k^2}}{2},\\
        \beta_k&=\dfrac{t_k-1}{t_{k+1}},\\
        y_k&=x_k+\beta_k(x_k-x_{k-1}),\\
        x_{k+1}&=\prox_{\tau g}\bigl(y_k-\tau\nabla f(y_k)\bigr).
    \end{cases}
\end{equation}
Jika \(g\equiv0\), metode ini sangat mirip dengan algoritma bola berat. Namun, gradien kini dievaluasi setelah suku inersia ditambahkan, dan konvergensinya bergantung pada pemilihan adaptif parameter inersia yang cermat.

% segment-id: d90.hab.v1.ch06.seg0009
\begin{lemma}
    Tetapkan \(t_1=1\). Maka
    \begin{equation}
        t_k\geq\frac{k+1}{2}
    \end{equation}
    untuk setiap \(k\geq1\).
\end{lemma}
\begin{proof}
    Bukti mengikuti induksi langsung. Pernyataan benar untuk \(k=1\). Dengan memakai hipotesis induksi, kita memperoleh
    \begin{equation}
        t_{k+1}
        =\frac{1+\sqrt{1+4t_k^2}}{2}
        \geq\frac{1+\sqrt{1+(k+1)^2}}{2}
        \geq\frac{k+2}{2}.
    \end{equation}
\end{proof}

% segment-id: d90.hab.v1.ch06.seg0010
\begin{lemma}[Ketaksamaan fundamental gradien proksimal]
    Definisikan
    \begin{equation}
        P_\tau(x)=\prox_{\tau g}\bigl(x-\tau\nabla f(x)\bigr).
    \end{equation}
    Di bawah hipotesis komposit di atas, untuk setiap \(x,y\in\R^d\) dan \(0<\tau\leq1/L\), berlaku
    \begin{equation}
        F(x)-F(P_\tau(y))
        \geq\frac{1}{2\tau}\|x-P_\tau(y)\|^2
        -\frac{1}{2\tau}\|x-y\|^2+\ell_f(x,y),
    \end{equation}
    dengan \(\ell_f(x,y)=f(x)-f(y)-\langle x-y,\nabla f(y)\rangle\).
\end{lemma}
\begin{proof}
    Pertama, kita dapat menulis
    \begin{equation}
        \begin{aligned}
            P_\tau(y)
            &=\arg\min_z
            \frac{1}{2\tau}\|z-(y-\tau\nabla f(y))\|^2+g(z)\\
            &=\arg\min_z
            f(y)+\langle z-y,\nabla f(y)\rangle
            +\frac{1}{2\tau}\|z-y\|^2+g(z).
        \end{aligned}
    \end{equation}
    Sebagai catatan, identitas ini menunjukkan bahwa metode gradien proksimal dapat ditafsirkan sebagai langkah proksimal terhadap \(g\) ditambah pendekatan linear terhadap \(f\). Definisikan
    \begin{equation}
        \phi(z)=f(y)+\langle z-y,\nabla f(y)\rangle
        +\frac{1}{2\tau}\|z-y\|^2+g(z).
    \end{equation}
    Karena \(g\) konveks, \(\phi\) bersifat \(1/\tau\)-konveks kuat. Oleh karena itu,
    \begin{equation}
        \phi(x)-\phi(P_\tau(y))
        \geq\frac{1}{2\tau}\|x-P_\tau(y)\|^2.
    \end{equation}
    Lebih lanjut, karena \(f\) konveks dan \(L\)-halus serta \(1/\tau\geq L\), lema penurunan memberikan
    \begin{equation}
        f(z)\leq f(y)+\langle z-y,\nabla f(y)\rangle
        +\frac{L}{2}\|z-y\|^2
        \leq f(y)+\langle z-y,\nabla f(y)\rangle
        +\frac{1}{2\tau}\|z-y\|^2.
    \end{equation}
    Khususnya, \(\phi(P_\tau(y))\geq F(P_\tau(y))\), sehingga
    \begin{equation}
        \phi(x)-F(P_\tau(y))
        \geq\frac{1}{2\tau}\|x-P_\tau(y)\|^2.
    \end{equation}
    Dengan memasukkan definisi \(\phi\) dan memakai
    \(f(y)+\langle x-y,\nabla f(y)\rangle=f(x)-\ell_f(x,y)\),
    penataan ulang hubungan terakhir tepat memberikan
    \begin{equation}
        F(x)-F(P_\tau(y))
        \geq\frac{1}{2\tau}\|x-P_\tau(y)\|^2
        -\frac{1}{2\tau}\|x-y\|^2+\ell_f(x,y),
    \end{equation}
    dan bukti selesai.
\end{proof}

% segment-id: d90.hab.v1.ch06.seg0011
\begin{theorem}[Konvergensi FISTA]
    Misalkan \(f:\R^d\rightarrow\R\) konveks, terdiferensialkan, dan bergradien \(L\)-Lipschitz dengan \(L>0\); misalkan \(g:\R^d\rightarrow(-\infty,\infty]\) proper, semikontinu bawah, dan konveks; serta ambil \(x^*\in\arg\min F\). Inisialisasikan FISTA dengan \(x_0=x_1\), \(t_1=1\), dan ukuran langkah tetap \(0<\tau\leq1/L\). Untuk setiap \(k\geq2\),
    \begin{equation}
        F(x_k)-F(x^*)\leq\frac{C}{\tau k^2},
        \qquad
        C=2\left(t_2^2\|x^*-x_1\|^2+
        \|t_2(x^*-x_2)+(t_2-1)(x_1-x^*)\|^2\right).
    \end{equation}
\end{theorem}

% segment-id: d90.hab.v1.ch06.seg0012
\begin{proof}
    Pilih \(x=t_{k+1}^{-1}x^*+(1-t_{k+1}^{-1})x_k\) dan \(y=y_k\) dalam ketaksamaan fundamental gradien proksimal. Karena \(f\) konveks, \(\ell_f(x,y)\geq0\). Maka
    \begin{equation}\label{eq:fista1}
        \begin{aligned}
            F(t_{k+1}^{-1}x^*+&(1-t_{k+1}^{-1})x_k)-F(x_{k+1})\\
            \geq{}&
            \frac{1}{2\tau}
            \|t_{k+1}^{-1}x^*+(1-t_{k+1}^{-1})x_k-x_{k+1}\|^2\\
            &-\frac{1}{2\tau}
            \|t_{k+1}^{-1}x^*+(1-t_{k+1}^{-1})x_k-y_k\|^2\\
            ={}&
            \frac{1}{2\tau t_{k+1}^2}
            \|x^*+(t_{k+1}-1)x_k-t_{k+1}x_{k+1}\|^2\\
            &-\frac{1}{2\tau t_{k+1}^2}
            \|x^*+(t_{k+1}-1)x_k-t_{k+1}y_k\|^2.
        \end{aligned}
    \end{equation}
    Di sisi lain, menurut kekonveksan \(F\) dan fakta \(t_k\geq1\),
    \begin{equation}\label{eq:fista2}
        \begin{aligned}
            F(t_{k+1}^{-1}x^*+(1-t_{k+1}^{-1})x_k)-F(x_{k+1})
            \leq{}&t_{k+1}^{-1}F(x^*)
            +(1-t_{k+1}^{-1})F(x_k)-F(x_{k+1})\\
            ={}&(1-t_{k+1}^{-1})(F(x_k)-F(x^*))
            -(F(x_{k+1})-F(x^*)).
        \end{aligned}
    \end{equation}
    Karena \(y_k=x_k+\frac{t_k-1}{t_{k+1}}(x_k-x_{k-1})\), kita mempunyai
    \begin{equation}\label{eq:fista3}
        \begin{aligned}
            \|x^*+(t_{k+1}-1)x_k-t_{k+1}y_k\|^2
            =\|x^*-t_kx_k+(t_k-1)x_{k-1}\|^2.
        \end{aligned}
    \end{equation}
    Menggabungkan~\eqref{eq:fista1}, \eqref{eq:fista2}, dan \eqref{eq:fista3} memberikan
    \begin{equation}
        \begin{aligned}
            t_{k+1}^2(1-t_{k+1}^{-1})
            &(F(x_k)-F(x^*))
            -t_{k+1}^2(F(x_{k+1})-F(x^*))\\
            \geq{}&
            \frac{1}{2\tau}
            \|x^*-t_{k+1}x_{k+1}+(t_{k+1}-1)x_k\|^2\\
            &-\frac{1}{2\tau}
            \|x^*-t_kx_k+(t_k-1)x_{k-1}\|^2.
        \end{aligned}
    \end{equation}
    Dari pembaruan FISTA berlaku \(t_{k+1}^2(1-t_{k+1}^{-1})=t_k^2\). Jadi,
    \begin{equation}
        \begin{aligned}
            t_k^2&(F(x_k)-F(x^*))
            -t_{k+1}^2(F(x_{k+1})-F(x^*))\\
            \geq{}&
            \frac{1}{2\tau}
            \|x^*-t_{k+1}x_{k+1}+(t_{k+1}-1)x_k\|^2\\
            &-\frac{1}{2\tau}
            \|x^*-t_kx_k+(t_k-1)x_{k-1}\|^2.
        \end{aligned}
    \end{equation}
    Dengan demikian,
    \begin{equation}
        \begin{aligned}
            t_k^2&(F(x_k)-F(x^*))
            +\frac{1}{2\tau}
            \|x^*-t_kx_k+(t_k-1)x_{k-1}\|^2\\
            \geq{}&
            t_{k+1}^2(F(x_{k+1})-F(x^*))
            +\frac{1}{2\tau}
            \|x^*-t_{k+1}x_{k+1}+(t_{k+1}-1)x_k\|^2.
        \end{aligned}
    \end{equation}
    Dengan mengiterasikan ketaksamaan ini, untuk \(k\geq2\) diperoleh
    \begin{equation}
        \begin{aligned}
            t_2^2&(F(x_2)-F(x^*))
            +\frac{1}{2\tau}
            \|t_2(x^*-x_2)+(t_2-1)(x_1-x^*)\|^2\\
            \geq{}&
            t_k^2(F(x_k)-F(x^*))
            +\frac{1}{2\tau}
            \|x^*-t_kx_k+(t_k-1)x_{k-1}\|^2.
        \end{aligned}
    \end{equation}
    Selain itu, ketaksamaan fundamental gradien proksimal dengan \(x=x^*\) dan \(y=y_1=x_0=x_1\) memberikan
    \begin{equation}
        F(x^*)-F(x_2)
        \geq\frac{1}{2\tau}\|x^*-x_2\|^2
        -\frac{1}{2\tau}\|x^*-x_1\|^2.
    \end{equation}
    Dengan membuang suku kuadrat nonnegatif pada ruas kanan estimasi energi dan memakai batas terakhir, kita menemukan
    \begin{equation}
        \begin{aligned}
            F(x_k)-F(x^*)
            \leq t_k^{-2}\left(
            \frac{t_2^2}{2\tau}\|x^*-x_1\|^2
            +\frac{1}{2\tau}
            \|t_2(x^*-x_2)+(t_2-1)(x_1-x^*)\|^2
            \right).
        \end{aligned}
    \end{equation}
    Karena \(t_k\geq(k+1)/2\), batas dalam pernyataan teorema mengikuti.
\end{proof}


\backmatter
\chapter*{Catatan koreksi terhadap sumber}
Terjemahan ini mempertahankan isi sumber kecuali pada koreksi yang dapat ditentukan berikut. Koreksi dilakukan oleh penyusun edisi bahasa Indonesia dan bukan perubahan yang disahkan oleh penulis sumber.

\begin{enumerate}
    \item Batas bawah kompleksitas pembuka kini memulihkan urutan kuantor bahwa untuk setiap metode terdapat fungsi, memuat norma kuadrat yang hilang, dan menyatakan secara eksplisit data awal serta hipotesis bahwa fungsi konveks, bergradien \(L\)-Lipschitz, mempunyai peminim dan dimensinya cukup besar untuk jumlah iterasi yang ditinjau.
    \item Bukti rumus Gelfand memakai norma operator yang jelas dan kompleksifikasi untuk nilai eigen nonreal. Bentuk Jordan hanya menempatkan satu pada superdiagonal di dalam satu blok, memakai indeks nilpotensi yang benar, menangani \(\rho(A)=0\) secara terpisah, dan memperbaiki pemilihan pangkat ekstrem pada batas.
    \item Bukti korolari radius spektral membuktikan arah \(A^k\rightarrow0\Rightarrow\rho(A)<1\) melalui vektor eigen dan menangani indeks awal dalam batas geometrik secara eksplisit.
    \item Teorema bola berat kini mensyaratkan fungsi \(C^2\) di sekitar peminim stasioner, \(0<\mu\leq L\), serta parameter konstan. Klaim parameter optimal dibatasi pada radius spektral kasus terburuk bagi model kuadratik atau linearisasi lokal.
    \item Rekursi keadaan tidak lagi memfaktorkan gradien nonlinear sebagai matriks. Selisih gradien dilinearkan melalui Hessian, residu didefinisikan dengan tepat, dan tanda gangguan keadaan diperbaiki menjadi \((-\tau r_k,0)\).
    \item Polinom karakteristik diturunkan tanpa membagi dengan nilai eigen, sehingga kasus \(\beta=0\) tercakup. Batas akar yang keliru diganti dengan kriteria Schur--Jury dan bukti kestabilan lokal dalam norma ekuivalen.
    \item Parameter kasus terburuk kini disertai bukti batas bawah minimaks yang lengkap melalui kriteria Schur--Jury terskala, termasuk kasus kelengkungan tunggal \(L=\mu\); kondisi kesamaan menghasilkan penyeimbangan kedua ujung spektrum \([\mu,L]\). Klaim optimalitas tidak lagi digeneralisasi menjadi jaminan global bagi semua fungsi konveks mulus.
    \item Bagian FISTA kini menyatakan domain, kekonveksan, semikontinuan bawah, keberadaan peminim, \(L>0\), \(0<\tau\leq1/L\), inisialisasi, dan rentang indeks secara lengkap.
    \item Bukti ketaksamaan fundamental memperbaiki frasa kekonveksan \(g\), memulihkan kuadrat yang hilang pada \(\|z-y\|^2\), serta menata ulang substitusi fungsi model dengan benar.
    \item Teorema FISTA kini menyatakan semua hipotesis komposit, menetapkan \(k\geq2\), dan memberi konstanta eksplisit yang bergantung pada inisialisasi. Identitas energi dan teleskop sumber dipertahankan.
    \item Salah ketik pada notasi waktu, frasa \emph{bergantung pada iterasi}, komponen inersia, aljabar linear, pernyataan korolari, kemiripan FISTA, dan teks bukti diperbaiki tanpa mengubah isi matematis lainnya.
\end{enumerate}

\chapter*{Atribusi dan lisensi}
Karya sumber: Andreas Habring, \emph{Lecture Notes: Convex Optimization}, arXiv:2607.11664v1, 13 Juli 2026, \url{https://arxiv.org/abs/2607.11664v1}. Karya sumber tersedia berdasarkan Creative Commons Attribution 4.0 International (CC BY 4.0), \url{https://creativecommons.org/licenses/by/4.0/}.

Edisi bahasa Indonesia ini merupakan karya turunan. Hak cipta materi sumber tetap pada pemegang haknya. Anda boleh membagikan dan mengadaptasinya dengan atribusi yang sesuai, tautan lisensi, dan penandaan perubahan sebagaimana disyaratkan CC BY 4.0. Catatan koreksi dan materi penghubung baru dinyatakan sebagai perubahan edisi; tidak ada dukungan atau pengesahan oleh penulis sumber maupun institusinya.

\end{document}
